% Report on Use of AI Tools
% Per COMAP 2025 Contest Policy - This section is outside the 25-page limit

\section*{Report on Use of AI Tools}
\addcontentsline{toc}{section}{Report on Use of AI Tools}

In accordance with COMAP's 2025 policy on AI tool usage, we disclose the following applications of artificial intelligence tools during the development of this solution.

\vspace{1em}

\noindent\textbf{1. Claude (Anthropic, Claude-3.5 Sonnet, January 2025 version)}

\noindent\textit{Purpose:} Literature review, mathematical model development, code implementation assistance, and writing refinement.

\noindent\textbf{Query 1:}
\begin{quote}
``Summarize the key methods for State of Charge (SOC) estimation in lithium-ion batteries, including Coulomb counting, Kalman filtering, and model-based approaches. Include relevant academic references.''
\end{quote}

\noindent\textbf{Output 1:} The AI provided a structured summary of SOC estimation methods including: (1) Coulomb counting with integration error accumulation issues, (2) OCV-based methods using voltage-SOC relationships, (3) Extended Kalman Filter approaches combining models with measurements, and (4) equivalent circuit model methods. References to Plett (2004, 2015), Hu et al. (2012), and Chen \& Rinc\'on-Mora (2006) were suggested. \textit{We verified all references against original sources and incorporated relevant methods into our model.}

\vspace{0.5em}

\noindent\textbf{Query 2:}
\begin{quote}
``Derive the state-space equations for a 2-RC equivalent circuit model coupled with thermal dynamics, suitable for Extended Kalman Filter implementation.''
\end{quote}

\noindent\textbf{Output 2:} The AI provided the coupled ODE system with state vector $[SOC, V_{RC1}, V_{RC2}, T]^T$, including the Coulomb counting equation, RC branch dynamics, and lumped thermal model. \textit{We independently verified the equations through dimensional analysis and adapted them to include SOH dynamics for our aging model.}

\vspace{0.5em}

\noindent\textbf{Query 3:}
\begin{quote}
``Implement a 4th-order Runge-Kutta solver in Python for the coupled SOC-thermal ODE system with adaptive time-stepping.''
\end{quote}

\noindent\textbf{Output 3:} The AI provided Python code using scipy.integrate.solve\_ivp with RK45 method. \textit{We modified the code for our specific model parameters and validated against analytical solutions for simplified cases.}

\vspace{0.5em}

\noindent\textbf{Query 4:}
\begin{quote}
``Explain the Sobol sensitivity analysis method and how to implement it for battery model parameter importance ranking.''
\end{quote}

\noindent\textbf{Output 4:} The AI explained variance-based global sensitivity with first-order and total-order indices, and suggested using SALib Python library. \textit{We implemented the method following Saltelli et al. (2008) and verified results with manual parameter sweeps.}

\vspace{1em}

\noindent\textbf{2. GitHub Copilot (January 2025 version)}

\noindent\textit{Purpose:} Auto-completions for Python/MATLAB code used in preparing numerical simulations and data visualization.

\noindent\textbf{Usage:} Code auto-completion suggestions were used for boilerplate code in plotting functions, file I/O operations, and array manipulations. All auto-completed code was reviewed and tested before inclusion.

\vspace{1em}

\noindent\textbf{3. Baidu Translate / DeepL (2025 version)}

\noindent\textit{Purpose:} Translation assistance for technical terminology verification.

\noindent\textbf{Usage:} Used to verify translations of technical battery terminology between Chinese and English to ensure accuracy in the final English report.

\vspace{1em}

\noindent\textbf{4. MATLAB (R2024a with Symbolic Math Toolbox)}

\noindent\textit{Purpose:} Symbolic computation and numerical simulation.

\noindent\textbf{Usage:} Built-in functions including ode45 solver, symbolic differentiation for Jacobian computation, and optimization toolbox for parameter fitting. These are standard computational tools with embedded AI-assisted features.

\vspace{1em}

\noindent\textbf{5. Python Libraries (NumPy, SciPy, SALib)}

\noindent\textit{Purpose:} Numerical computation and sensitivity analysis.

\noindent\textbf{Usage:} Standard scientific computing libraries used for ODE solving, optimization, and Sobol sensitivity analysis implementation.

\vspace{1.5em}

\subsection*{Verification and Validation Statement}

All AI-assisted outputs were subject to rigorous verification:
\begin{itemize}[noitemsep]
    \item Mathematical equations were verified through dimensional analysis and boundary condition checks
    \item All literature citations were traced to original sources and verified for accuracy
    \item Numerical code was validated against analytical solutions where available
    \item Model predictions were compared against experimental data from published datasets
    \item Final writing was reviewed by all team members for technical accuracy
\end{itemize}

\subsection*{Statement of Originality}

The core intellectual contributions of this work---including problem interpretation, model architecture design, result analysis, and scientific conclusions---represent the original work of the team members. AI tools served as productivity aids; all strategic decisions and critical evaluations were performed by the human team.
